\section*{Overview}
This document outlines the revision strategy for the paper "DeepPrep: An LLM-Powered Agentic System for Autonomous Data Preparation". Based on the feedback from three reviewers (Weak Accept, Weak Reject, Weak Accept), the primary focus of the revision will be on **validating the necessity of a specialized agent** compared to general coding agents, **demonstrating robustness on real-world data**, and **analyzing scalability**.

\section{Action Items Checklist}

\subsection*{A. New Experiments}
\begin{itemize}
    \item \textbf{[Exp-1] Comparison with General Coding Agents (R1W1, R2W1):} 
    \begin{itemize}
        \item Implement or configure a baseline using a general-purpose agent framework (e.g., OpenDevin, MetaGPT, or a standard ReAct loop with GPT-4-Turbo/Claude-3.5-Sonnet).
        \item Evaluate on the existing benchmark to demonstrate why a specialized ADP agent (DeepPrep) is necessary (e.g., lower cost, higher success rate in handling schema constraints).
    \end{itemize}
    
    \item \textbf{[Exp-2] Real-World/Messy Data Evaluation (R1W3, R2W2, R3W4):}
    \begin{itemize}
        \item Select 10-20 complex, real-world data preparation scenarios (e.g., from Kaggle, WDC Web Tables, or messy Excel sheets).
        \item These scenarios should include ambiguity, missing values, and non-standard formatting not present in the synthetic SQL-derived datasets.
        \item Conduct a case study or a small-scale quantitative evaluation on these tasks.
    \end{itemize}
    
    \item \textbf{[Exp-3] Scalability Analysis (R1D3, R2W4, R3D3):}
    \begin{itemize}
        \item \textbf{Data Scalability:} Measure inference time and success rate as table size increases (1k $\to$ 10k $\to$ 100k rows).
        \item \textbf{Search Scalability:} Analyze how the tree grows (number of nodes visited) relative to the pipeline length.
    \end{itemize}
    
    \item \textbf{[Exp-4] Baseline Fairness (R2W3):}
    \begin{itemize}
        \item Perform a quick grid search or sensitivity check on baseline prompts to ensure they are not "strawman" implementations. Document this in the appendix.
    \end{itemize}
\end{itemize}

\subsection*{B. Content Revisions}
\begin{itemize}
    \item \textbf{[Text-1] Technical Clarifications (R1D1, R1D2, R2D1, R3D1):}
    \begin{itemize}
        \item Explicitly detail the prompt structure: Schema + Sample Rows (not full content).
        \item Explain the tree search strategy (e.g., DFS/BFS, pruning limits).
        \item Clarify the "probabilistic" formulation vs. deterministic execution.
    \end{itemize}
    
    \item \textbf{[Text-2] Novelty Argumentation (R1W1, R2W1, R3W1):}
    \begin{itemize}
        \item Rewrite the introduction and related work to clearly distinguish DeepPrep from Tree-of-Thoughts (ToT) and general agents.
        \item Highlight: Execution-grounded state representation (materialized intermediate tables) vs. pure text reasoning in ToT.
    \end{itemize}
    
    \item \textbf{[Text-3] Failure Analysis (R1D5):}
    \begin{itemize}
        \item Add a subsection discussing typical failure modes (e.g., semantic misunderstanding of column names vs. syntax errors).
    \end{itemize}
\end{itemize}

\newpage

\section{Detailed Responses to Reviewers}

\subsection*{Response to Reviewer \#1 (Weak Accept)}

\noindent \textbf{R1W1: Comparison against general-purpose coding agents (e.g., Claude Code).} \\
\textit{Response Strategy:} We acknowledge this gap. We will add a comparison against a strong general-purpose agent baseline (e.g., an agent loop using GPT-4o or Claude-3.5-Sonnet with full code execution capabilities). We hypothesize that while general agents are capable, DeepPrep's specialized state representation (schema-aware) and constrained operator space result in higher reliability and lower token costs for data prep tasks.

\noindent \textbf{R1W2 \& D1: Missing technical details (prompts, state representation, schema handling).} \\
\textit{Response Strategy:} We will add a dedicated "Implementation Details" section in the Appendix.
\begin{itemize}
    \item \textbf{Prompts:} We will publish the exact prompt templates.
    \item \textbf{Content vs. Schema:} We will clarify that the LLM sees the \textit{schema} and \textit{first $k$ rows} (samples), not the full table, to manage context window limits.
    \item \textbf{Intermediate States:} We will clarify that intermediate schemas are derived by actually executing the generated code on the sample data (execution-grounded).
\end{itemize}

\noindent \textbf{R1D2: Benefit of tree structure and reuse of failed branches.} \\
\textit{Response Strategy:} We will clarify that the tree structure allows for \textit{non-local backtracking}. Unlike linear retry loops, if a path proves dead (e.g., leads to an empty table or schema mismatch), the agent can revert to a parent node (a previous valid table state) and try a different operator. We will explicitly state whether failed nodes are added to the context as negative constraints (if applicable) or simply pruned.

\noindent \textbf{R1D3: Sensitivity and Scalability (Turns, Table Size).} \\
\textit{Response Strategy:} See \textbf{[Exp-3]}. We will add a "Scalability" subsection in the experiments. We will plot "Success Rate vs. Max Turns" and "Time vs. Row Count". We expect the system to be robust to row count (since it operates on samples/logic) but sensitive to column count (context length).

\noindent \textbf{R1D5: Failure Analysis.} \\
\textit{Response Strategy:} We will add a qualitative analysis of failure cases, categorizing them into: (1) Ambiguous User Intent, (2) Hallucinated Column Names, and (3) Logic Errors that produce valid but incorrect SQL/Pandas code.

\subsection*{Response to Reviewer \#2 (Weak Reject)}

\noindent \textbf{R2W1 \& D1: Limited conceptual novelty; resembles existing frameworks.} \\
\textit{Response Strategy:} We will sharpen our novelty claims. While the components (LLM, ReAct, Tree Search) exist, the \textit{integration} for Data Prep is novel. Specifically, we will highlight:
\begin{itemize}
    \item Unlike generic code agents, our "Environment" is the Dataframe itself.
    \item Our "State" is not just text history, but the materialized intermediate data schema.
    \item The progressive training curriculum is specific to the complexity of data transformation chains.
\end{itemize}

\noindent \textbf{R2W2 \& D2: Reliance on synthetic datasets; need real-world data.} \\
\textit{Response Strategy:} This is the most critical critique. See \textbf{[Exp-2]}. We cannot re-do the entire training on real data due to label scarcity, but we will add a robust \textit{evaluation} on a set of real-world "messy" datasets (e.g., WDC, TPC-DI, or Kaggle) to prove the model generalizes beyond the synthetic distribution.

\noindent \textbf{R2W3 \& D3: Baseline prompt tuning details.} \\
\textit{Response Strategy:} We will expand the experimental setup section. We will explicitly state the prompt engineering efforts for baselines (e.g., "We used Chain-of-Thought prompting for all baselines and optimized the few-shot examples"). We will provide the baseline prompts in the supplementary material to ensure reproducibility.

\noindent \textbf{R2W4 \& D4: Scalability analysis (Data size, Pipeline length).} \\
\textit{Response Strategy:} Addressed in \textbf{[Exp-3]}. We will provide graphs showing the system's behavior as the pipeline length increases (demonstrating the benefit of the tree search over linear prompting for long-horizon tasks).

\subsection*{Response to Reviewer \#3 (Weak Accept)}

\noindent \textbf{R3W1: Difference from Tree-of-Thoughts (ToT) / MCTS.} \\
\textit{Response Strategy:} We will add a discussion in the Related Work.
\begin{itemize}
    \item \textbf{ToT:} Usually performs search in the \textit{reasoning space} (text thoughts) without external feedback.
    \item \textbf{DeepPrep:} Performs search in the \textit{action space} (code execution) with \textit{grounded feedback} (execution results/errors). The nodes in our tree are concrete data states, not just abstract thoughts.
\end{itemize}

\noindent \textbf{R3W2: Probabilistic formula vs. Deterministic execution.} \\
\textit{Response Strategy:} We will clarify the formalism. The \textit{policy} (Agent) is stochastic (probabilistic generation of code). The \textit{transition} (Environment execution) is deterministic given the code and data. The formulation intends to model the agent's uncertainty in selecting the correct operator, not randomness in the SQL engine.

\noindent \textbf{R3W3: Scalability and Tree Pruning.} \\
\textit{Response Strategy:} We will add details on the search control. We likely use a heuristic (e.g., beam width or max depth) to prevent explosion. We will explicitly state these hyperparameters (e.g., "Max depth = 5, Beam width = 3") and discuss the computational complexity $O(Depth \times Beam)$.

\noindent \textbf{R3W4: More evaluation on independent real-world data.} \\
\textit{Response Strategy:} Addressed in \textbf{[Exp-2]}. We will emphasize that the new real-world case studies are "out-of-domain" relative to our training data.
